\chapter{Diffusion Models: Intuition}

\section{Introduction}

Diffusion models generate data by reversing a gradual corruption process. Instead of producing a sample in a single step, they start from noise and iteratively refine it into a structured output.

This chapter develops an intuitive understanding of diffusion models. We explain progressive noising, denoising, and the idea of score-based learning, and connect diffusion to iterative refinement.

\section{Why Generate by Denoising?}

Directly generating complex data (e.g., images) in one step is difficult. Diffusion models take a different approach:
\begin{itemize}
    \item start from simple noise,
    \item gradually transform it into structured data,
    \item learn each small transformation.
\end{itemize}

This breaks a difficult problem into many simpler steps.

\section{Forward Diffusion as Gradual Corruption}

The forward process adds noise step by step:
\[
x_0 \to x_1 \to x_2 \to \cdots \to x_T.
\]

Each step adds small Gaussian noise:
\[
x_t = \sqrt{1 - \beta_t} \, x_{t-1} + \sqrt{\beta_t} \, \epsilon_t,
\quad \epsilon_t \sim \mathcal{N}(0,I).
\]

\subsection{Interpretation}

\begin{itemize}
    \item early steps: slightly corrupted data,
    \item later steps: heavily corrupted,
    \item final state: nearly pure Gaussian noise.
\end{itemize}

Thus the forward process gradually destroys structure.

\subsection{Gaussian Corruption Sequence (Example)}

Starting from an image:
\begin{itemize}
    \item $x_0$: clean image,
    \item $x_1$: slightly noisy,
    \item $x_{10}$: visibly corrupted,
    \item $x_T$: indistinguishable from noise.
\end{itemize}

\section{Reverse Denoising as Learned Reconstruction}

The generative process reverses this corruption:
\[
x_T \to x_{T-1} \to \cdots \to x_0.
\]

A neural network learns to remove noise at each step:
\[
p_\theta(x_{t-1} \mid x_t).
\]

\subsection{Interpretation}

Each step performs:
\begin{itemize}
    \item a small denoising operation,
    \item a refinement of structure,
    \item a move toward the data distribution.
\end{itemize}

Thus generation is iterative reconstruction.

\section{Noise Prediction and Score Intuition}

Instead of directly predicting $x_{t-1}$, diffusion models often predict the noise $\epsilon$ added at each step.

\subsection{Noise Prediction}

Given $x_t$, the model predicts $\epsilon_t$ such that:
\[
x_0 \approx \frac{x_t - \sqrt{\beta_t}\,\epsilon_t}{\sqrt{1-\beta_t}}.
\]

\subsection{Score Interpretation}

The model implicitly learns the \emph{score function}:
\[
\nabla_x \log p(x),
\]
which points toward regions of higher probability.

\begin{itemize}
    \item denoising corresponds to moving toward higher-density regions,
    \item each step nudges samples toward realistic configurations.
\end{itemize}

Thus diffusion can be viewed as gradient ascent in probability space with noise.

\section{Iterative Refinement View}

Diffusion models generate by repeated small corrections:
\[
x_{t-1} = x_t + \text{small update}.
\]

This resembles:
\begin{itemize}
    \item residual refinement,
    \item gradient-based optimization,
    \item iterative reconstruction algorithms.
\end{itemize}

\subsection{Thought Experiment}

Imagine starting from random pixels:
\begin{itemize}
    \item each step removes a bit of noise,
    \item patterns begin to emerge,
    \item structure gradually sharpens,
    \item final result resembles real data.
\end{itemize}

This incremental process is easier to learn than direct generation.

\section{Why Diffusion Models Work Well in Practice}

Several factors contribute to their success:

\begin{itemize}
    \item \textbf{Stable training:} objective resembles supervised denoising,
    \item \textbf{Well-behaved gradients:} avoids adversarial instability,
    \item \textbf{Iterative refinement:} decomposes complex generation,
    \item \textbf{Flexibility:} can model diverse data distributions.
\end{itemize}

Compared to GANs:
\begin{itemize}
    \item more stable,
    \item less prone to mode collapse,
    \item but slower at sampling.
\end{itemize}

\section{Worked Example: Small-Step Denoising}

Suppose:
\[
x_t = x_0 + \epsilon,
\quad \epsilon \sim \mathcal{N}(0, \sigma^2 I).
\]

If $\sigma$ is small:
\begin{itemize}
    \item $x_t$ is close to $x_0$,
    \item predicting $\epsilon$ is easy,
    \item recovering $x_0$ is accurate.
\end{itemize}

Thus learning small denoising steps is easier than large jumps.

\section{Comparison with Other Generative Models}

\begin{itemize}
    \item \textbf{Autoencoders:} reconstruct inputs in one step,
    \item \textbf{GANs:} generate samples via adversarial training,
    \item \textbf{Diffusion:} generate via many small denoising steps.
\end{itemize}

Diffusion trades speed for stability and quality.

\section{A Unifying Perspective}

Diffusion models view generation as reversing noise. The forward process simplifies data into noise; the reverse process reconstructs structure step by step.

This connects to:
\begin{itemize}
    \item score matching (learning gradients of log-density),
    \item iterative optimization,
    \item dynamical systems.
\end{itemize}

Thus diffusion provides a conceptually simple and powerful generative paradigm.

\section{Summary}

In this chapter we developed an intuitive understanding of diffusion models. We described progressive noising, reverse denoising, and noise prediction. We connected these ideas to score functions and iterative refinement.

Diffusion models succeed by decomposing generation into many small, learnable steps, providing stability and high-quality outputs.

\section*{Exercises}

\begin{enumerate}
    \item Explain why small-step denoising is easier to learn than direct generation.

    \item Consider a Gaussian corruption process. Show how variance increases over steps.

    \item Compare diffusion models and autoencoders in terms of reconstruction.

    \item Compare diffusion models and GANs in terms of:
    \begin{enumerate}
        \item training stability,
        \item sample quality,
        \item computational cost.
    \end{enumerate}

    \item Describe how iterative refinement improves generation quality.

    \item Explain the role of noise prediction in diffusion models.

    \item Give an intuitive explanation of the score function $\nabla_x \log p(x)$.

    \item Why do diffusion models avoid mode collapse more easily than GANs?
\end{enumerate}